% =============================================================
% 6Q Proposal — Contact-centred benchmark for soft vs rigid
% grippers on thin, flat cloth.
%
% Source: 6Q_Grasping cloth.pdf (project proposal, six-question
% structure). LaTeX port with citations to Clark et al. 2023
% (Dyson Household Clothing Set) and Teeple et al. 2022 (Harvard
% Multi-Dimensional Compliance) + comparison table.
% =============================================================

\documentclass[10pt,letterpaper]{article}

\usepackage[margin=1in]{geometry}
\usepackage{cite}
\usepackage{amsmath,amssymb,amsfonts}
\usepackage{graphicx}
\usepackage{textcomp}
\usepackage{xcolor}
\usepackage{booktabs}
\usepackage{array}
\usepackage{tabularx}
\usepackage{enumitem}
\usepackage{caption}
\usepackage[hidelinks]{hyperref}

\newcolumntype{L}[1]{>{\raggedright\arraybackslash}p{#1}}

\title{A Contact-Centred Benchmark for Soft Passive-Adaptive vs Rigid
       Parallel-Jaw Grippers on Thin, Flat Cloth}

\author{Shiyao Ni \\
        McGill University \\
        \texttt{shiyao.ni@mail.mcgill.ca}}
\date{}

\begin{document}
\maketitle

\begin{abstract}
Robotic manipulation of thin, flat cloth resting on a hard surface remains
an open problem. Most existing solutions augment rigid parallel-jaw
grippers with environment-assisted actions, specialised end-effectors, or
closed-loop perception, treating the gripper itself as a fixed component.
Prior benchmarks for cloth manipulation either standardise tasks across
many gripper designs without instrumenting the contact interaction
\cite{Clark2023}, or compare soft and rigid grippers under a single
fixed-condition test with an actively-controlled pneumatic soft hand
\cite{Teeple2022}. This proposal targets the contact stage directly:
a benchmark that explicitly treats gripper compliance as the independent
variable and maps the per-gripper success envelope in a five-axis
control-parameter space, comparing a soft, passively-adaptive UMI
finray (TPU~90A) against a factory Franka parallel jaw on the same
Franka Research~3 arm under an identical control policy. The output is
a pair of envelopes whose shape and area form the basis of the soft-vs-rigid
comparison, complemented by three functional cloth-manipulation tasks
under a shared human-defined policy.
\end{abstract}

\section{Problem}
\label{sec:problem}
Current robot-manipulation research relies heavily on the operator to hand
the target object into the manipulator before any trajectory is executed.
This work targets the upstream step that is typically skipped: the robot
picking up a thin, flat cloth lying on a hard surface, without
human staging and without active replanning. The specific failure mode
under study is the contact stage of the grasp, where the gripper must
make controlled contact with the cloth--surface system and pinch a
graspable layer.

\section{Motivation}
\label{sec:motivation}
If a robot cannot pick up cloth flat from a surface, the practical
autonomy of downstream manipulation pipelines is severely constrained.
Successful flat-cloth grasping enables:
\begin{itemize}[leftmargin=*,nosep]
    \item Laundry folding and textile handling;
    \item Hospital bedding management;
    \item Assistive robotics for elderly or disabled individuals;
    \item Garment manufacturing and logistics;
    \item Soft-object packaging.
\end{itemize}
Each application sits behind the same gating capability: lifting a
single flat layer from a hard surface reliably.

\section{State of the Art}
\label{sec:sota}
Current solutions combine environment interaction, specialised
end-effectors, and feedback mechanisms to make cloth grasping
tractable.

\textbf{Environment-assisted actions} create graspable features
(folds, wrinkles, exposed edges) through dragging, wiping, or peeling,
allowing a standard parallel-jaw gripper to perform a pinch on a
prepared edge.

\textbf{Specialised end-effectors} reduce reliance on precise edge access:
variable-friction grippers, micro-needle or needle-based graspers,
electroadhesion, and airflow-based methods (Bernoulli or Coand\u{a}
effects).

\textbf{Closed-loop perception and control} verify grasp success
(e.g., single-ply detection) and trigger recovery actions when failures
occur. In practice, most systems still rely on rigid parallel-jaw
grippers augmented with these strategies.

Two recent works are directly relevant to this proposal because they
attempt cross-gripper comparison rather than perception or planning
improvements. Clark~\textit{et al.}~\cite{Clark2023} introduce a
standardised household-clothing object set and four benchmark tasks
(edge drag accuracy, edge grasp resilience, crumpled object
encapsulation, flat non-boundary grasp success), and evaluate four
parallel grippers including a Fin~Ray\textregistered{} variant. However,
their benchmarks score outcomes (placement offset, payload, droop,
binary success) rather than instrument the contact stage; the soft-vs-rigid
difference is reported but not isolated from the task-specific
gripper kinematics.

Teeple~\textit{et al.}~\cite{Teeple2022} directly compare a rigid
parallel jaw to a soft pneumatic two-finger gripper, showing that
vertical, lateral, and rotational compliance each protect against
distinct failure modes (positional uncertainty, snag duration, peak
snag force). Their soft gripper, however, is actively pressure-controlled,
and the comparison is performed at a fixed vertical sweep without a
multi-axis control-parameter envelope.

\section{Drawbacks of the State of the Art}
\label{sec:drawbacks}
\begin{itemize}[leftmargin=*,nosep]
    \item Heavy reliance on environment-assisted actions reduces true
          autonomy and inflates task complexity and execution time.
    \item Specialised end-effectors are task-specific, hard to
          generalise, and often fragile or impractical for real-world
          deployment.
    \item Rigid parallel-jaw grippers require precise positioning and
          well-defined geometry, making them unreliable for thin, flat
          cloth on a hard surface.
    \item Closed-loop verification and recovery introduce sensing and
          control overhead.
    \item Most approaches treat the gripper as fixed and focus on
          perception or planning, so there is no systematic evaluation
          of how gripper design (rigid vs.\ soft) affects the contact
          stage in isolation.
\end{itemize}
Consequently, directly grasping thin, flat cloth from a surface remains
an open problem. The work that does compare soft and rigid grippers
either does not instrument the contact \cite{Clark2023} or relies on
active control of the soft side \cite{Teeple2022}, leaving the
contribution of a \emph{passive} soft gripper at the contact stage
uncharacterised.

\section{Contribution}
\label{sec:contribution}
This proposal contributes a benchmark focused on the contact stage of
grasping thin, flat cloth, comparing a soft but passively-adaptive
gripper (UMI finray, TPU~90A; no added control space relative to a
parallel jaw) against a rigid parallel-jaw gripper on the same robot
arm under the same human-defined control policy.

\subsection{Benchmark Metrics}
\textit{Control-policy admissibility:} a trial is admissible only if
the Franka controller did not quit due to a force-safety error during
execution. Per-trial outcome is determined by a researcher visual
verification of grasp success, combined with the controller-status flag.

The benchmark sweeps five gripper-control axes (approach z-offset,
head rotation, applied-force setpoint, finger gap, closing speed) and
records the boundary of the region in which both conditions hold. Five
tolerance metrics per gripper are then read off the envelope:
\begin{enumerate}[leftmargin=*,nosep]
    \item \textbf{z-axis tolerance:} range of approach-height offset
          over which the grasp succeeds.
    \item \textbf{Angular tolerance:} orientation spread of the
          Franka achieved end-effector pose observed across successful
          trials with commanded orientation held constant (a passive
          observation, not a swept parameter).
    \item \textbf{Maximum admissible surface-contact force} at which the
          controller still agrees (does not quit).
    \item \textbf{Peak contact force} during successful grasps (table or
          cloth side).
    \item \textbf{Minimum contact force} required for a successful grasp.
\end{enumerate}

\subsection{Functional Tasks}
Both grippers are additionally tested on three downstream tasks under
a single shared human-defined policy:
\begin{itemize}[leftmargin=*,nosep]
    \item Diagonal fold (flat cloth corner grasp);
    \item Middle pick $\rightarrow$ drop onto adjacent sphere (flat
          cloth middle grasp);
    \item Drag from A to B (flat cloth edge grasp).
\end{itemize}
A future scope extension may add tasks dealing with folded cloth.

\subsection{Comparison with Prior Benchmarks}
Table~\ref{tab:comparison} situates this work alongside the two most
directly comparable prior benchmarks. Conceptually, the proposal sits
closer to Teeple~\textit{et al.}~\cite{Teeple2022} (same soft-vs-rigid
question on flat thin cloth) than to Clark~\textit{et al.}~\cite{Clark2023}
(broad standardised benchmark across many gripper designs), but it
replaces Teeple's two-point comparison with a five-axis swept envelope
and instruments the under-cloth force-distribution side with a
distributed force-sensor array rather than a single wrist
force/torque sensor.

Distinctive contributions of this work relative to both prior benchmarks:

\begin{enumerate}[leftmargin=*,nosep]
    \item \emph{Envelope shape, not aggregate pass-rate.} The output is
          a region in 5-D control-parameter space whose boundary is the
          metric, not a per-task scalar score.
    \item \emph{Controller-quit as a tracked admissibility gate.}
          Neither prior benchmark explicitly separates ``controller
          refused to continue'' from ``grasp failed.'' Teeple's
          rigid-gripper failures include UR5e safety-stop trips that
          are conflated with grasp failure. This benchmark separates
          them to avoid biasing the comparison against grippers
          operating near the controller's force-safety threshold.
    \item \emph{Distributed under-cloth force sensing.} A five-channel
          force array (Tekscan A301-1) measures per-foot table-side
          reaction force; a two-channel A301-25 pair measures
          finger-pad pinch force; the Franka built-in $F_{\text{ext}}$
          and joint torques are logged in parallel.
    \item \emph{Force traceability.} Per-channel calibration is
          traceable to an OIML M1 mass standard via dead-weight loading
          with a power-law fit; calibration freshness is a hard gate
          on data admissibility.
    \item \emph{Passive soft gripper.} The soft side is a passively
          deforming UMI finray, not an active pneumatic system,
          isolating the contribution of compliance to grasp success
          without introducing additional control variables.
\end{enumerate}

\begin{table}[!t]
\centering
\caption{Comparison of this proposal with two directly relevant prior benchmarks for cloth grasping.}
\label{tab:comparison}
\footnotesize
\renewcommand{\arraystretch}{1.25}
\begin{tabularx}{\textwidth}{@{}L{2.5cm} X X X@{}}
\toprule
\textbf{Axis} & \textbf{This proposal} & \textbf{Clark et al. 2023~\cite{Clark2023}} & \textbf{Teeple et al. 2022~\cite{Teeple2022}} \\
\midrule
Research question
  & Map the control-parameter success envelope per gripper (size and shape of the region in which the grasp succeeds and the Franka controller does not quit).
  & Define a standardised cloth-manipulation benchmark suite (object set + edge taxonomy + four tasks) for cross-paper gripper comparison.
  & Show that vertical, lateral, and rotational compliance each protect against distinct failure modes (uncertainty, snag duration, peak snag force). \\
Output type
  & Two parameter-space envelopes, one per gripper, compared in shape and area.
  & Per-gripper score on each of four tasks (offset~mm, payload~N, droop~mm, success~\%).
  & Range of vertical-offset tolerance (mm); snag displacement (mm, deg), force duration (s), peak lateral force (N). \\
Object set
  & 2 cloth specimens (1 thin + 1 thick), flat only.
  & 13-item household clothing set with edge classification (single/double/folded; boundary/non-boundary).
  & 3 swatches: elastic spandex, woven cotton, folded woven cotton. \\
Cloth state(s) tested
  & Flat only (crumpled/folded out of scope).
  & Flat boundary edges, non-boundary edges, crumpled, compressed-crumpled.
  & Flat only; snags imposed by external clamp. \\
Grippers compared
  & 22~mm UMI finray (TPU~90A), 17~mm UMI finray (TPU~90A), Franka factory parallel jaw (\textbf{3}).
  & OpenHand T42, Robotiq 2F-140, FE~Hand, FE~Hand (Fin~Ray Flex) (\textbf{4}).
  & Rigid PhantomX AX-12 (foam pads) vs custom soft pneumatic 2-finger (\textbf{2}). \\
Independent variable
  & \emph{Swept:} approach z-offset, head rotation, applied force, finger gap, closing speed.
  & \emph{Categorical:} object $\times$ edge type $\times$ gripper.
  & \emph{Swept:} vertical centring offset (\(-40\) to \(+4\)~mm); other parameters fixed. \\
Robot
  & Franka Research~3.
  & Franka Emika Panda.
  & UR5e (150~N safety stop). \\
Surface / fixturing
  & Instrumented 3D-printed test bench; rigid feet on A301-1 array (TPU~95A 1~mm contact pad option per Bill~0003 amendment).
  & A0 alignment grid on rigid backing plate.
  & Plain tabletop. \\
Force instrumentation
  & 5$\times$ A301-1 table feet + 2$\times$ A301-25 finger pads + Franka $F_{\text{ext}}$ + Franka joint torques, shared ROS timebase; 100~Hz DAQ + 1~kHz Franka FCI.
  & Single load cell on motorised rail (resilience benchmark only); 1~mm ruler for placement.
  & UR5e wrist F/T sensor (500~Hz) + AprilTag pose tracking (30~Hz). \\
Force calibration rigor
  & OIML~M1 traceable dead-weight per-channel power-law fit; JSON record with derivation header; 30-day re-cal cadence; $\sigma$-discard threshold.
  & Not specified.
  & Stiffness measured ($n{=}3$); per-trial force traceability not specified. \\
Pass/fail criterion
  & Researcher visual verification + Franka controller-not-quit admissibility gate (treated as separate failure class).
  & Binary success per task.
  & Binary ``arm picks up swatch.'' \\
Distinctive failure modes tracked
  & Controller-quit / force-safety abort tracked as a signal-derived admissibility class.
  & Drop during motion; lift failure; encapsulation failure.
  & UR5e safety-stop and grasp instability (conflated with grasp failure for the rigid gripper); tear under snag. \\
Trials per condition
  & Parameter sweep (count per axis TBD at execution).
  & 5 repeats per object.
  & $n{=}3$ per condition. \\
\bottomrule
\end{tabularx}
\end{table}

\section{How the Contribution Addresses the Drawbacks}
\label{sec:addresses}
\begin{itemize}[leftmargin=*,nosep]
    \item \textbf{Gripper as a variable.} By explicitly treating the
          gripper as the independent variable rather than a fixed
          subsystem, this work addresses the central limitation
          identified in Section~\ref{sec:drawbacks}: that gripper
          design is rarely the swept axis in cloth-manipulation
          benchmarks.
    \item \textbf{Quantified behaviour under uncertainty.} The
          proposed metrics (positional tolerance, angular tolerance,
          force profiles) directly quantify how each gripper behaves
          under control-parameter uncertainty. This reduces reliance
          on environment-assisted actions such as dragging or folding
          to expose graspable features.
    \item \textbf{Isolated compliance effect.} By comparing a soft
          passively-adaptive gripper to a rigid parallel jaw under
          the same control policy, the study isolates the
          contribution of compliance and contact mechanics, rather
          than conflating it with perception or planning improvements.
          This is the gap left open by~\cite{Clark2023} (does not
          instrument contact) and~\cite{Teeple2022} (active pneumatic
          control on the soft side).
    \item \textbf{Force-related failure modes.} The minimum, peak, and
          maximum-admissible force metrics provide visibility into
          slipping, snagging, and excessive-force failure modes that
          are not surfaced by current benchmarks.
    \item \textbf{Task-level usefulness.} The three functional tasks
          ensure that grasps are not only successful at contact but
          also useful for downstream manipulation, addressing the gap
          between grasp acquisition and task execution.
    \item \textbf{Hardware-aware design.} The overall framing shifts
          part of the cloth-manipulation problem from complex
          perception and control toward hardware-aware design,
          potentially reducing system complexity while improving
          robustness on thin, deformable objects.
\end{itemize}

\section{Grasping Strategy}
\label{sec:strategy}
The proposed contact protocol is a multi-stage descent followed by a
slow close and slow ascent. Defining the cloth--surface contact plane
as $z = 0$:
\begin{enumerate}[leftmargin=*,nosep]
    \item \textbf{Multi-stage descent.} Descend rapidly to the $(x, y)$
          contact position at $z = +5$~mm. Slow descent to $z = 0$~mm.
          Very slow descent to a small negative height (e.g.,
          $z = -3$~mm). This provides a safety margin and a force
          limit, and allows the finray to deform slightly. The
          resulting contact force between the lower section of the
          gripper and the cloth enables the subsequent steps.
    \item \textbf{Slow close.} Close the gripper slowly. Because
          contact force already exists between the lower section of
          the gripper and the cloth, this closing motion pinches the
          cloth into the inner contact surface of the gripper.
    \item \textbf{Slow ascent.} Ascend slowly in the first few steps to
          allow the finray to recover its original shape and grip the
          cloth firmly.
    \item \textbf{Pinch margin.} The margin between the gripper's
          current opening and the cloth thickness determines the
          amount of fabric pinched into the contact zone.
    \item \textbf{Comparison with parallel jaw.} The protocol is
          enabled by the deformable nature of the finray. The Franka
          parallel jaw can succeed on the same task but does not
          require (and cannot exploit) the negative-height step.
          Without that step, the parallel jaw requires careful
          height tuning: too-high fails on insufficient contact;
          too-low fails on excessive contact force triggering the
          controller-quit admissibility gate. Quantifying this
          tuning margin per gripper is precisely the z-axis tolerance
          metric in Section~\ref{sec:contribution}.
\end{enumerate}

\begin{thebibliography}{9}

\bibitem{Clark2023}
A.~B. Clark, L.~Cramphorn-Neal, M.~Rachowiecki, and A.~Gregg-Smith,
``Household clothing set and benchmarks for characterising end-effector
cloth manipulation,'' in \emph{Proc.\ IEEE Int.\ Conf.\ Robot.\ Autom.\
(ICRA)}, 2023.

\bibitem{Teeple2022}
C.~B. Teeple, J.~Werfel, and R.~J. Wood,
``Multi-dimensional compliance of soft grippers enables gentle
interaction with thin, flexible objects,'' in \emph{Proc.\ IEEE Int.\
Conf.\ Robot.\ Autom.\ (ICRA)}, 2022, pp.~728--734.

\end{thebibliography}

\end{document}
